\section*{Capítulo \ref{ch:b2:energy-geology} --- Geologia da energia}
% numbers from the Book 4 second-writer scratch calculator calc/ch30.py

\begin{solution}{exo:b2:energy-geology:1}
(a) 1.58: tipo~I; 0.95: tipo~III; 1.24: tipo~II; 1.12: tipo misto~II/III. O querogênio de tipo~III, com
0.95, produz sobretudo gás. (b) H/C $= 150/100 = 1.50$, O/C $= 8/100 = 0.08$. (c) Por \qty{100}{g}:
$77.4/12.011 = \qty{6.444}{mol}$ de C, $9.8/1.008 = \qty{9.722}{mol}$ de H, $8.2/15.999 =
\qty{0.513}{mol}$ de O; H/C $= 1.51$ e O/C $= 0.0795$, de acordo com a fórmula.
\end{solution}

\begin{solution}{exo:b2:energy-geology:2}
(a) 0.45: imaturo, abaixo do início em \qty{0.60}{\percent}; 0.68: gerando óleo; 1.1: ainda na
janela do petróleo, perto de seu fim; 1.6: além de 1.20--\qty{1.30}{\percent}, de modo que o óleo se decompõe em metano e a
rocha produz gás; 2.5: só gás, com a geração prosseguindo até cerca de \qty{4.8}{\percent}. (b) Não: a geração
de metano pelo tipo~III começa perto de \qty{0.73}{\percent}. (c) $F = (\ln R_o + 1.6)/3.7$: $(-0.511 +
1.6)/3.7 = 0.294$ e $(0.262 + 1.6)/3.7 = 0.503$.
\end{solution}

\begin{solution}{exo:b2:energy-geology:3}
(a) No reservatório, $6 \times 10^6 \times 25 \times 0.18 \times 0.70 = \qty{1.89e7}{m^3}$;
$\mathrm{STOIIP} = 1.89 \times 10^7/1.25 = \qty{1.51e7}{m^3}$. (b) $0.30 \times 1.512 \times 10^7 =
\qty{4.54e6}{m^3}$ e $0.40 \times 1.512 \times 10^7 = \qty{6.05e6}{m^3}$.
\end{solution}

\begin{solution}{exo:b2:energy-geology:4}
(a) $V = 10^7 \times 200 = \qty{2e9}{m^3}$ e $q_R = 2.5 \times 10^6 \times 2 \times 10^9 \times 60
= \qty{3.0e17}{J}$. (b) Abaixo de \qty{90}{\celsius}: um \omterm{def:b2:energy-geology:geothermal}{recurso de baixa entalpia}, para uso direto do
calor. (c) $P = 0.25 \times 3.0 \times 10^{17}/(30 \times \qty{3.156e7}{s}) = \qty{79.2}{MW}$ de calor.
\end{solution}

\begin{solution}{exo:b2:energy-geology:5}
(a) $k_2/k_1 = \exp[(E/R)(1/T_1 - 1/T_2)] = 2$ com $1/373.15 - 1/383.15 = \qty{6.994e-5}{K^{-1}}$
dá $E = 8.314 \times \ln 2/(6.994 \times 10^{-5}) = \qty{82.4}{kJ/mol}$. (b) $\exp(26800 \times
6.994 \times 10^{-5}) = 6.52$. (c) A regra da duplicação, cuja energia efetiva é mais baixa, dá
relativamente mais peso aos longos tempos passados a baixa temperatura; para um querogênio de tipo~I de alta
energia de ativação, o TTI anuncia a geração cedo demais.
\end{solution}

\begin{solution}{exo:b2:energy-geology:6}
(a) $x = 1/2$ na aproximação dá $(ART^2/\beta E)\,e^{-E/RT} = \ln 2$; tomando logaritmos,
$E/RT = \ln[AT^2/(\beta(E/R)\ln 2)]$, donde a forma enunciada. Partindo de qualquer estimativa perto de
\qty{400}{K} e iterando com $\beta$ em kelvins por segundo: a \qty{1}{\celsius\per\Myr}
($\beta = \qty{3.169e-14}{K/s}$), $T = \qty{409.6}{K}$, \qty{136.4}{\celsius}; a \qty{10}{K} por
minuto ($\beta = \qty{0.1667}{K/s}$, $\num{5.26e12}$ vezes mais rápido), $T = \qty{719.0}{K}$,
\qty{445.9}{\celsius}. (b) A diferença é de \qty{309.5}{K}; a estimativa $6.259 \times
\ln(5.26 \times 10^{12}) = 6.259 \times 29.29 = \qty{183}{K}$ é pequena demais, porque o próprio $RT^2/E$
cresce com $T$ (até \qty{19.3}{K} a \qty{719}{K}): o deslocamento linear só vale para pequenas
variações de taxa.
\end{solution}

\begin{solution}{exo:b2:energy-geology:7}
$\lambda_5 - \lambda_8 = \qty{8.29725e-10}{yr^{-1}}$ e $r_0 = 0.7204/99.2742 = 0.0072567$.
(a) Há \qty{2.0}{\Ga}, $r = 0.0072567\,e^{1.6595} = 0.03814$, proporção $0.03814/1.03814 =
\qty{3.67}{\percent}$; há \qty{1.8}{\Ga}, $r = 0.03231$, proporção \qty{3.13}{\percent}. (b) $r =
3/97 = 0.030928$ e $t = \ln(0.030928/0.0072567)/(8.29725 \times 10^{-10}) = \qty{1.75}{\Ga}$.
(c) O minério empobrecido e o urânio normal decaíram desde então pelos mesmos fatores, de modo que a razão entre
suas proporções de $^{235}$U é hoje a mesma que logo após as reações: o minério tinha então metade da
proporção normal, de modo que cerca de metade de seu $^{235}$U havia sofrido fissão.
\end{solution}

\begin{solution}{exo:b2:energy-geology:8}
(a) $10^9/1.872 = \qty{5.34e8}{m^3}$; $10^9/621.4 = \qty{1.61e6}{m^3}$; $10^9/686.4 =
\qty{1.46e6}{m^3}$. (b) $7 \times 10^9/583 = \qty{1.20e7}{m^3}$ de \ce{CO2}, que preenchem
$0.35 \times 0.30 = 0.105$ do volume de rocha: $1.20 \times 10^7/0.105 = \qty{1.14e8}{m^3}$ de rocha,
uma espessura média de $1.143 \times 10^8/(5 \times 10^6) = \qty{22.9}{m}$.
\end{solution}

\begin{solution}{exo:b2:energy-geology:9}
(a) $V = \qty{9e9}{m^3}$ e $q_R = 2.6 \times 10^6 \times 9 \times 10^9 \times 235 =
\qty{5.50e18}{J}$. (b) Em $30 \times \qty{3.156e7}{s} = \qty{9.47e8}{s}$: $R_g = 0.1$ dá
\qty{581}{MW} e $0.25$ dá \qty{1450}{MW} de calor. (c) $6 \times (10 \text{ a } 20) = 60$ a
\qty{120}{MW} de eletricidade, $60/581 = 0.10$ a $120/581 = 0.21$ da potência térmica: a maior parte do
calor é rejeitada, como em qualquer usina que trabalha com uma diferença de temperatura modesta.
\end{solution}

\begin{solution}{exo:b2:energy-geology:10}
(a) A resistência é $0.6 \times (80 - 39) = \qty{24.6}{MPa} > \qty{22}{MPa}$: estável. A ruptura exige
$0.6\,\Delta p = 24.6 - 22$, $\Delta p = \qty{4.33}{MPa} = \qty{43}{bar}$, da mesma ordem que os
\qty{32}{bar} que delimitaram os sismos de Denver. (b) $(24.6 - 24)/0.6 = \qty{1.0}{MPa} =
\qty{10}{bar}$: uma falha já perto da ruptura precisa só de um pequeno aumento. (c) O fluido injetado eleva
a pressão numa zona que se propaga por difusão através da rocha permeável e das zonas de falha, a quilômetros
de distância ao longo de anos; onde quer que encontre uma falha já perto da ruptura, um aumento modesto basta.
\end{solution}

\begin{solution}{exo:b2:energy-geology:11}
(a) Escreva $x = e^{-u}$ com $u = (ART^2/\beta E)\,e^{-E/RT}$. Uma taxa de transformação de
\qty{10}{\percent} significa $x = 0.9$, $u = \ln(10/9)$; \qty{90}{\percent} significa $u = \ln 10$. Como
$\mathrm{d}\ln u/\mathrm{d}T = E/RT^2 + 2/T \approx E/RT^2$, $\Delta(\ln u) \approx (E/RT^2)\,\Delta T$,
o que dá a fórmula, com $\ln\ln 10 - \ln\ln(10/9) = 0.834 + 2.250 = 3.084$. (b) Com $T =
\qty{409.6}{K}$, $RT^2/E = \qty{6.26}{K}$ e $\Delta T = \qty{19.3}{K}$; a curva exata de energia única
dá \qty{18.6}{K}, de 125.3 a \qty{143.8}{\celsius}. (c) A curva contínua sobe de
\qty{10}{\percent} perto de \qty{122}{\celsius} a \qty{90}{\percent} perto de \qty{148}{\celsius}, ao longo de cerca de
\qty{26}{K}: a dispersão das energias de ativação alarga o intervalo de conversão.
\end{solution}

\begin{solution}{exo:b2:energy-geology:12}
(a) $2 \times 10^{10}/20 = \qty{1.0e9}{m^3}$ (\qty{1.0}{km^3}) e $2 \times 10^{10}/60 =
\qty{3.3e8}{m^3}$. (b) $2 \times 10^{10}/621.4 = \qty{3.22e7}{m^3}$. (c) $3.22 \times 10^7/0.02 =
\qty{1.61e9}{m^3}$ de poros, o que, com $50 \times 0.30 = \qty{15}{m}$ de espaço poroso por metro
quadrado, exige $1.61 \times 10^9/15 = \qty{1.07e8}{m^2}$, cerca de \qty{107}{km^2}.
\end{solution}

\begin{solution}{pb:b2:energy-geology:1}
\noindent\textbf{Parte I.}
\textbf{1.} A \qty{70}{\Ma}: $50 \times 60 = \qty{3000}{m}$, $15 + 25 \times 3.0 =
\qty{90}{\celsius}$. Hoje: $3000 + 70 \times 30 = \qty{5100}{m}$, \qty{142.5}{\celsius}.
\textbf{2.} $60 \times 0.025 = \qty{1.5}{\celsius\per\Myr}$, depois $30 \times 0.025 =
\qty{0.75}{\celsius\per\Myr}$.
\textbf{3.} 15--\qty{20}{\celsius}: \qty{3.33}{\Myr}; cada intervalo de 20 a \qty{90}{\celsius}:
$10/1.5 = \qty{6.67}{\Myr}$; cada um de 90 a \qty{140}{\celsius}: $10/0.75 = \qty{13.3}{\Myr}$;
140--\qty{142.5}{\celsius}: \qty{3.33}{\Myr}.
\textbf{4.} Valores por intervalo $\Delta t_n 2^n$, com as somas acumuladas: 15--20 ($n = -9$) 0.0065; 20--30
0.026 (0.033); 30--40 0.052 (0.085); 40--50 0.104 (0.189); 50--60 0.208 (0.397); 60--70 0.417
(0.814); 70--80 0.833 (1.65); 80--90 1.67 (3.31); 90--100 6.67 (9.98); 100--110 13.3 (23.3);
110--120 26.7 (50.0); 120--130 53.3 (103); 130--140 107 (210); 140--142.5 ($n = 4$) 53.3: TTI
atual 263.
\textbf{5.} O valor 15 é ultrapassado no intervalo 100--\qty{110}{\celsius}, a $100 + 10 \times
(15 - 9.98)/13.33 = \qty{103.8}{\celsius}$, atingidos $(103.8 - 90)/0.75 = \qty{18.4}{\Myr}$ depois de
\qty{70}{\Ma}: a \qty{51.6}{\Ma} e a $(103.8 - 15)/0.025 = \qty{3550}{m}$.
\textbf{6.} 75: \qty{124.7}{\celsius}, \qty{23.7}{\Ma}, \qty{4390}{m}. 160: \qty{135.3}{\celsius},
\qty{9.6}{\Ma}, \qty{4810}{m}.
\textbf{7.} Não: com 263 ela já passou do fim da geração de óleo (160), desde cerca de \qty{10}{\Ma}.

\noindent\textbf{Parte II.}
\textbf{8.} $k = 9.21 \times 10^{13}\,e^{-26800/363.15} = \qty{8.20e-19}{s^{-1}}$. A rocha esteve
abaixo de \qty{90}{\celsius} durante os \qty{50}{\Myr} anteriores a \qty{70}{\Ma}, de modo que a fração convertida é
menor que $k\,t = 8.20 \times 10^{-19} \times 50 \times 3.156 \times 10^{13} = \num{1.29e-3}$.
\textbf{9.} $\beta = 0.75/(3.156 \times 10^{13}) = \qty{2.376e-14}{K/s}$, $T = \qty{415.65}{K}$:
$ART^2/(\beta E) = \num{2.498e28}$ e $e^{-E/RT} = \num{9.95e-29}$, logo $u = 2.486$, $x = 0.0833$ e
a taxa de transformação é 0.917.
\textbf{10.} \qty{134.7}{\celsius}.
\textbf{11.} $70 - (134.7 - 90)/0.75 = \qty{10.4}{\Ma}$.
\textbf{12.} O TTI situa o pico de geração a \qty{23.7}{\Ma}, a cinética \qty{13}{\Myr} mais tarde,
porque a regra da duplicação corresponde a $E \approx \qty{82}{kJ/mol}$ e pondera demais as baixas
temperaturas para um querogênio de \qty{223}{kJ/mol} (exercício~\ref{exo:b2:energy-geology:5}).
\textbf{13.} Mais cedo: uma energia mais baixa abaixa a temperatura de meia conversão.

\noindent\textbf{Parte III.}
\textbf{14.} Rocha: $4 \times 10^8 \times 50 \times 2400 = \qty{4.8e13}{kg}$; potencial $0.030
\times 4.8 \times 10^{13} = \qty{1.44e12}{kg}$; gerado $0.92 \times 1.44 \times 10^{12} =
\qty{1.32e12}{kg}$.
\textbf{15.} $1.3248 \times 10^{12}/850 = \qty{1.56e9}{m^3}$.
\textbf{16.} $1.2 \times 10^7 \times 40 \times 0.22 \times 0.75/1.30 = \qty{6.09e7}{m^3}$.
\textbf{17.} $6.09 \times 10^7/1.56 \times 10^9 = 0.039$, \qty{3.9}{\percent}: a maior parte do óleo gerado
fica na rocha geradora, perde-se ao longo dos caminhos de migração ou escapa para a superfície.
\textbf{18.} $0.35 \times 6.09 \times 10^7 = \qty{2.13e7}{m^3}$.
\textbf{19.} Cerca de dois terços do óleo ficam no reservatório após a recuperação primária e a
secundária; a injeção de \ce{CO2} poderia recuperar mais 10 a \qty{15}{\percent} do óleo in situ.

\noindent\textbf{Parte IV.}
\textbf{20.} $15 + 25 \times 2 = \qty{65}{\celsius}$ e $0.101 + 1000 \times 9.807 \times 2000/10^6
= \qty{19.7}{MPa}$: ambos acima do ponto crítico (\qty{31.0}{\celsius}, \qty{7.38}{MPa}), de modo que o
\ce{CO2} é supercrítico, a cerca de \qty{686}{kg.m^{-3}}.
\textbf{21.} $2.13 \times 10^7 \times 1.30 = \qty{2.77e7}{m^3}$.
\textbf{22.} Volume poroso $1.2 \times 10^7 \times 40 \times 0.22 = \qty{1.056e8}{m^3}$, e $E =
2.77 \times 10^7/1.056 \times 10^8 = 0.262$, muito acima dos 1 a \qty{4}{\percent} das formações
salinas: o óleo produzido esvaziou um espaço conhecido e vedado.
\textbf{23.} $6 \times 10^8 \times 80 \times 0.20 \times 691 \times 0.02 = \qty{1.33e11}{kg}$,
\qty{133}{Mt}.
\textbf{24.} A \qty{1}{Mt} por ano, cerca de 133 anos.
\textbf{25.} $M = \rho\,N_p\,B_o = 686 \times 2.77 \times 10^7 = \qty{1.90e10}{kg}$: o reservatório
depletado pode armazenar \textbf{\qty[text-series-to-math=true]{19.0}{Mt}} de \ce{CO2}.
\end{solution}
